\section{Root space decomposition}

\begin{definition}
    A Lie algebra $L$ is called \emph{toral} if $\operatorname{ad}(a)$ is semisimple
    for every $a\in L$.
\end{definition}

\begin{lemma}
    Any toral Lie algebra is abelian.
\end{lemma}

\begin{proof}
    Fix $0 \neq a \in L$. Since $\operatorname{ad} a$ is semisimple,
    \[
        L=\bigoplus_{\lambda} L_\lambda(a),
        \qquad
        L_\lambda(a)=\{x\in L\mid [a,x]=\lambda x\}.
    \]
    We claim that $L_\alpha(a)=0$ for all $\alpha\neq 0$.

    Suppose instead that $L_\alpha(a)\neq 0$ for some $\alpha\neq 0$, and choose
    $0\neq b\in L_\alpha(a)$. If $x\in L_\beta(a)$, then the Jacobi identity gives
    \[
        [a,[b,x]]
        =
        [[a,b],x]+[b,[a,x]]
        =
        \alpha[b,x]+\beta[b,x]
        =
        (\alpha+\beta)[b,x].
    \]
    Hence
    \[
        [b,L_\beta(a)]\subseteq L_{\alpha+\beta}(a).
    \]
    Thus $\operatorname{ad} b$ carries each eigenspace of $\operatorname{ad} a$ into the
    next one shifted by $\alpha$. Since $\operatorname{ad} a$ has only finitely many
    eigenvalues, it follows that $\operatorname{ad} b$ is nilpotent.

    On the other hand, $L$ is toral, so $\operatorname{ad} b$ is semisimple. Hence
    $\operatorname{ad} b=0$. But
    \[
        (\operatorname{ad} b)(a)=[b,a]=-\alpha b\neq 0,
    \]
    a contradiction. Therefore $L_\alpha(a)=0$ whenever $\alpha\neq 0$.

    So $L=L_0(a)$, i.e.
    \[
        [a,x]=0
        \qquad\text{for all }x\in L.
    \]
    Since $a$ was arbitrary, $L$ is abelian.
\end{proof}

Now let $L$ be semisimple, and let $H$ be a maximal toral subalgebra of $L$.
Since $H$ is toral, it is abelian by the lemma above. Hence the commuting
semisimple endomorphisms $\operatorname{ad}(h)$ $(h\in H)$ are simultaneously
diagonalizable, and we obtain a decomposition
\[
    L=L_0\oplus \bigoplus_{\alpha\in \Delta} L_\alpha,
\]
where
\[
    L_\alpha=\{x\in L\mid [h,x]=\alpha(h)x \text{ for all } h\in H\},
    \qquad \alpha\in H^*,
\]
and $\Delta\subseteq H^*\setminus\{0\}$ is the set of nonzero weights.


\begin{lemma}
    For $\alpha,\beta\in \Delta\cup\{0\}$ one has
    \[
        [L_\alpha,L_\beta]\subseteq L_{\alpha+\beta},
    \]
    where by convention $L_{\alpha+\beta}=0$ if $\alpha+\beta\notin \Delta\cup\{0\}$.
\end{lemma}

\begin{proof}
    Take $x_\alpha\in L_\alpha$, $y_\beta\in L_\beta$, and $h\in H$. Then
    \[
        [h,[x_\alpha,y_\beta]]
        =
        [[h,x_\alpha],y_\beta]+[x_\alpha,[h,y_\beta]]
        =
        \alpha(h)[x_\alpha,y_\beta]+\beta(h)[x_\alpha,y_\beta].
    \]
    So
    \[
        [h,[x_\alpha,y_\beta]]
        =
        (\alpha+\beta)(h)[x_\alpha,y_\beta].
    \]
    Hence $[x_\alpha,y_\beta]\in L_{\alpha+\beta}$.
\end{proof}

Let $(\cdot\mid\cdot)$ denote the Killing form of $L$.

\begin{lemma}
    For any $\alpha,\beta\in \Delta\cup\{0\}$, if $\alpha+\beta\neq 0$, then
    \[
        (L_\alpha\mid L_\beta)=0.
    \]
\end{lemma}

\begin{proof}
    Choose $h\in H$ such that $(\alpha+\beta)(h)\neq 0$.
    Let $x_\alpha\in L_\alpha$ and $y_\beta\in L_\beta$. Then
    \[
        \alpha(h)(x_\alpha\mid y_\beta)
        =
        ([h,x_\alpha]\mid y_\beta)
        =
        (h\mid [x_\alpha,y_\beta]),
    \]
    while
    \[
        \beta(h)(x_\alpha\mid y_\beta)
        =
        (x_\alpha\mid [h,y_\beta])
        =
        -([h,x_\alpha]\mid y_\beta).
    \]
    Hence
    \[
        (\alpha(h)+\beta(h))(x_\alpha\mid y_\beta)=0.
    \]
    Since $(\alpha+\beta)(h)\neq 0$, we obtain
    \[
        (x_\alpha\mid y_\beta)=0.
    \]
\end{proof}

\begin{lemma}
    \[
        L_0=H.
    \]
\end{lemma}

\begin{proof}
    First note that the restriction of the Killing form to $L_0$ is nondegenerate.
    Indeed, if $x\in L_0$ and $(x\mid L_0)=0$, then by the previous lemma
    \[
        (x\mid L_\alpha)=0
        \qquad (\alpha\neq 0),
    \]
    so $x$ is orthogonal to all of $L$. Since the Killing form on $L$ is nondegenerate,
    $x=0$.

    Now let $a\in L_0$, and write its abstract Jordan decomposition in $L$ as
    \[
        a=a_s+a_n.
    \]
    Because $\operatorname{ad}(a)$ commutes with every $\operatorname{ad}(h)$ $(h\in H)$,
    the same is true of its semisimple and nilpotent parts. Hence
    \[
        a_s,a_n\in L_0.
    \]
    Since $a_s$ is semisimple and centralizes $H$, the subalgebra
    \[
        H+Fa_s
    \]
    is toral. By maximality of $H$, we must have $a_s\in H$.

    We claim that $a_n=0$. To prove this, it is enough --- by the nondegeneracy
    of $(\cdot\mid\cdot)$ on $L_0$ --- to show that
    \[
        (a_n\mid L_0)=0.
    \]

    First observe that $L_0$ is solvable. Indeed, if $u\in [L_0,L_0]$, then
    $u\in L_0$, so its semisimple part $u_s$ lies in $H$. On the other hand, since
    $u=[x,y]$ with $x,y\in L_0$ and $H$ centralizes $L_0$, one sees from the
    Killing-form computation that $u_s$ acts trivially on $L$; hence $u_s=0$.
    Therefore every element of $[L_0,L_0]$ is nilpotent, so Engel's theorem
    implies that $[L_0,L_0]$ is nilpotent. Thus $L_0$ is solvable.

    By Lie's theorem, there exists a basis of $L$ with respect to which all operators
    $\operatorname{ad}(x)$ $(x\in L_0)$ are upper triangular. Since $a_n$ is nilpotent,
    $\operatorname{ad}(a_n)$ is strictly upper triangular. Therefore, for every
    $x\in L_0$,
    \[
        (a_n\mid x)
        =
        \operatorname{tr}\bigl(\operatorname{ad}(a_n)\operatorname{ad}(x)\bigr)=0.
    \]
    Hence $(a_n\mid L_0)=0$, and so $a_n=0$.

    Thus $a=a_s\in H$. Since $a\in L_0$ was arbitrary, we have $L_0\subseteq H$.
    The reverse inclusion is clear, so $L_0=H$.
\end{proof}

\begin{corollary}
    The restriction of the Killing form to $H$ is nondegenerate.
\end{corollary}

\begin{proof}
    This follows immediately from the lemma $L_0=H$ and the nondegeneracy
    of the Killing form on $L_0$.
\end{proof}

\begin{corollary}
    For every $\alpha\in\Delta$, the pairing
    \[
        L_\alpha\times L_{-\alpha}\to F,
        \qquad (x,y)\mapsto (x\mid y),
    \]
    is nondegenerate. In particular, $L_{-\alpha}\neq 0$.
\end{corollary}

\begin{proof}
    Let $0\neq x\in L_\alpha$ and assume $(x\mid L_{-\alpha})=0$. By the orthogonality
    lemma, $x$ is also orthogonal to $H=L_0$ and to every $L_\beta$ with
    $\beta\neq -\alpha$. Hence $x$ is orthogonal to all of $L$, contradicting the
    nondegeneracy of the Killing form.
\end{proof}

We may therefore identify $H$ with $H^*$ via the Killing form:
\[
    \nu\colon H\longrightarrow H^*,
    \qquad
    h\longmapsto (h\mid \cdot)\vert_H .
\]
Thus
\[
    \nu(h)(y)=(h\mid y)
    \qquad \text{for all } y\in H.
\]
Let
\[
    \nu^{-1}\colon H^*\longrightarrow H
\]
be its inverse. For $\alpha,\beta\in H^*$ define
\[
    (\alpha\mid\beta):=\bigl(\nu^{-1}(\alpha)\mid \nu^{-1}(\beta)\bigr).
\]

\begin{lemma}
    For $x_\alpha\in L_\alpha$ and $y_{-\alpha}\in L_{-\alpha}$,
    \[
        [x_\alpha,y_{-\alpha}]
        =
        (x_\alpha\mid y_{-\alpha})\,\nu^{-1}(\alpha).
    \]
\end{lemma}

\begin{proof}
    Since $[L_\alpha,L_{-\alpha}]\subseteq L_0=H$, both sides lie in $H$.
    Let $h\in H$. Then
    \[
        ([h,x_\alpha]\mid y_{-\alpha})
        =
        (h\mid [x_\alpha,y_{-\alpha}]).
    \]
    But
    \[
        [h,x_\alpha]=\alpha(h)x_\alpha,
    \]
    so
    \[
        \alpha(h)(x_\alpha\mid y_{-\alpha})
        =
        (h\mid [x_\alpha,y_{-\alpha}]).
    \]
    By definition of $\nu^{-1}(\alpha)$,
    \[
        \alpha(h)
        =
        (\nu^{-1}(\alpha)\mid h).
    \]
    Hence
    \[
        (h\mid [x_\alpha,y_{-\alpha}])
        =
        \bigl(h\mid (x_\alpha\mid y_{-\alpha})\,\nu^{-1}(\alpha)\bigr).
    \]
    Since the Killing form on $H$ is nondegenerate, we conclude that
    \[
        [x_\alpha,y_{-\alpha}]
        =
        (x_\alpha\mid y_{-\alpha})\,\nu^{-1}(\alpha).
    \]
\end{proof}

\begin{lemma}
    For every $\alpha\in\Delta$,
    \[
        (\alpha\mid\alpha)\neq 0.
    \]
\end{lemma}

\begin{proof}
    Choose $0\neq x_\alpha\in L_\alpha$. By the preceding corollary, there exists
    $y_{-\alpha}\in L_{-\alpha}$ such that
    \[
        (x_\alpha\mid y_{-\alpha})\neq 0.
    \]
    Hence, by the previous lemma,
    \[
        [x_\alpha,y_{-\alpha}]
        =
        (x_\alpha\mid y_{-\alpha})\,\nu^{-1}(\alpha)\neq 0.
    \]

    Suppose $(\alpha\mid\alpha)=0$. Then
    \[
        \alpha\bigl(\nu^{-1}(\alpha)\bigr)
        =
        \bigl(\nu^{-1}(\alpha)\mid \nu^{-1}(\alpha)\bigr)
        =
        (\alpha\mid\alpha)
        =
        0.
    \]
    Set
    \[
        L' := Fx_\alpha \oplus Fy_{-\alpha}\oplus F\nu^{-1}(\alpha).
    \]
    Then
    \[
        [\nu^{-1}(\alpha),x_\alpha]
        =
        \alpha(\nu^{-1}(\alpha))x_\alpha
        =
        0,
    \]
    and similarly
    \[
        [\nu^{-1}(\alpha),y_{-\alpha}]=0.
    \]
    Also
    \[
        [x_\alpha,y_{-\alpha}]
        \in F\nu^{-1}(\alpha).
    \]
    Thus $L'$ is a solvable subalgebra.

    Consider the adjoint representation of $L'$ on $L$. By Lie's theorem,
    the operators $\operatorname{ad}(L')$ are simultaneously upper triangularizable.
    Since
    \[
        \nu^{-1}(\alpha)\in [L',L'],
    \]
    its adjoint action is strictly upper triangular, hence nilpotent.
    But $\nu^{-1}(\alpha)\in H$, so $\operatorname{ad}(\nu^{-1}(\alpha))$ is semisimple.
    Therefore it must be $0$. This implies $\nu^{-1}(\alpha)=0$, hence $\alpha=0$,
    a contradiction.

    Therefore $(\alpha\mid\alpha)\neq 0$.
\end{proof}

\begin{proposition}
    For every root $\alpha\in\Delta$, one may choose root vectors
    \[
        x_\alpha\in L_\alpha,
        \qquad
        y_{-\alpha}\in L_{-\alpha},
    \]
    such that if
    \[
        h_\alpha := \frac{2}{(\alpha\mid\alpha)}\,\nu^{-1}(\alpha),
    \]
    then
    \[
        [x_\alpha,y_{-\alpha}]=h_\alpha,
        \qquad
        [h_\alpha,x_\alpha]=2x_\alpha,
        \qquad
        [h_\alpha,y_{-\alpha}]=-2y_{-\alpha}.
    \]
\end{proposition}

\begin{proof}
    Choose $0\neq x_\alpha\in L_\alpha$. By the nondegeneracy of the pairing
    $L_\alpha\times L_{-\alpha}$, choose $y_{-\alpha}\in L_{-\alpha}$ so that
    \[
        (x_\alpha\mid y_{-\alpha})=\frac{2}{(\alpha\mid\alpha)}.
    \]
    Then the preceding lemma gives
    \[
        [x_\alpha,y_{-\alpha}]
        =
        \frac{2}{(\alpha\mid\alpha)}\,\nu^{-1}(\alpha)
        =
        h_\alpha.
    \]

    Now
    \[
        [h_\alpha,x_\alpha]
        =
        \alpha(h_\alpha)x_\alpha
        =
        \frac{2}{(\alpha\mid\alpha)}
        \alpha(\nu^{-1}(\alpha))x_\alpha
        =
        \frac{2}{(\alpha\mid\alpha)}(\alpha\mid\alpha)x_\alpha
        =
        2x_\alpha.
    \]
    Similarly,
    \[
        [h_\alpha,y_{-\alpha}]
        =
        -\alpha(h_\alpha)y_{-\alpha}
        =
        -2y_{-\alpha}.
    \]
\end{proof}

